\section{视觉恢复与聚合细节}
\label{sec:supp-queries}

\paragraph{输入约定与预处理。}

输入是同步图像 $\mathbf I_{t,v}$、时间戳、内参 $\mathbf K_v$ 及相机到传感器支架的标定 $\mathbf T^r_{c_v}$。支架位姿 $\mathbf T^w_r(t)$ 在受控评价中统一给定，并以第一帧为参考；相机位姿为 $\mathbf T^w_{c_v}(t)=\mathbf T^w_r(t)\mathbf T^r_{c_v}$。变换作用于列向量，距离单位为米。图像采用校正后的针孔模型与线性 RGB；畸变、滚动快门和曝光差异在证据聚合之前处理。

\paragraph{共享视觉编码。}

Driving Log Encoder 使用图像 patch 划分~\cite{dosovitskiy2021vit}与 Transformer 注意力~\cite{vaswani2017attention}，并加入图像位置、相机标定和时间编码。图像内部注意力提取空间特征，第二个 Transformer 在不同视角与时刻的池化图像摘要之间交互，再将上下文送回 patch 特征。几何头预测相机光轴深度、法线、置信度、语义和实例归属。置信度决定哪些预测可以提供表面证据。

可学习的资产查询读取全部图像 token，形成轨迹级身份槽；每个槽再读取单个时刻的特征以估计状态。由此，持久身份、尺寸和人体体型与逐帧变化的根位姿、关节旋转分开表示。根位姿先在支架坐标中预测，再通过 $\mathbf T^w_r(t)$ 映射到世界坐标。人体体型码选择对应身份的模板，模板包含固定规范支撑、骨架拓扑和绑定。身份槽容量限定可表示的资产数量。

\paragraph{查询支撑。}

每个查询存储规范位置 $\mathbf q_i$、单元尺度 $\boldsymbol\delta_i$、类型 $k_i$、局部注意力组 $g_i$ 和参考世界坐标 $\mathbf q_i^w$。可靠深度样本提供表面锚点。背景锚点与有限世界网格结合；刚体锚点经逆 SE(3) 变换对齐后，与局部包围体中的网格结合。人体查询位于规范模板表面，天空查询为按立方体贴图面分组的单位方向。这些支撑为已见与缺失内容提供明确位置。

补全过程中布局保持固定。我们将查询放到第一帧的配置中以得到 $\mathbf q_i^w$，其中人体包含蒙皮变换。规范坐标标识组件内部的内容，参考世界坐标则使生成器能够联系该组件与周围场景。

\paragraph{运动感知证据聚合。}

对于每份观测，我们将规范查询放到对应的资产状态：

\begin{equation}
\mathbf p_i(t)=
\begin{cases}
\mathbf q_i,& k_i=\mathrm{background},\\
\operatorname{xyz}(\mathbf T^w_a(t)\bar{\mathbf q}_i),& k_i=\mathrm{rigid},\\
\operatorname{xyz}(\mathbf T^w_h(t)\mathbf A_i(t)\bar{\mathbf q}_i),& k_i=\mathrm{human}.
\end{cases}
\label{eq:query-placement}
\end{equation}

其中，$\bar{\mathbf q}_i$ 为齐次坐标，$\operatorname{xyz}$ 提取欧氏坐标，$\mathbf A_i(t)$ 为加权蒙皮变换。投影到相机后，得到深度 $z_{i,t,v}$ 以及采样特征、深度和置信度 $(\boldsymbol\phi_{i,t,v},d_{i,t,v},c_{i,t,v})$。给定容差 $\eta=\eta_{\rm abs}+\eta_{\rm rel}d$，定义 $r=(d-z)/\eta$。表面观测需位于图像内、深度有效、置信度足够，与查询的实例身份一致，并满足 $|r|\leq1$。背景查询还要求静态背景语义。令 $a_{i,t,v}$ 表示这一有效性判断，则聚合为

\begin{align}
w_{i,t,v}&=a_{i,t,v}c_{i,t,v}\exp(-r_{i,t,v}^2/2),\\
\mathbf h_i&=\frac{\sum_{t,v}w_{i,t,v}[\boldsymbol\phi_{i,t,v};r_{i,t,v}]}
{\max(\varepsilon,\sum_{t,v}w_{i,t,v})}.
\label{eq:aggregation}
\end{align}


查询置信度取被接受观测中的最大值。重建实验采用中心射线投票，深度容差为 $\eta=0.15+0.02d$ m，置信度阈值为 0.8。可靠查询中心射线满足下式时提供自由空间票：
\begin{equation}
z_{i,t,v}+\tfrac12\|\boldsymbol\delta_i\|_2<d_{i,t,v}-\eta.
\label{eq:free-space}
\end{equation}
观测权重为正的查询标为已见；其余查询至少获得两票时标为自由空间，否则保持未知。天空投影使用相机旋转、可靠天空语义和零深度残差，不接收自由空间票。

\paragraph{基于投影足迹的自由空间审计。}
在 100 个合成测试场景的固定查询布局上，比较中心射线投票与足迹判据。足迹判据依据查询的刚体或蒙皮变换放置单元八个顶点，并计算
\begin{equation}
z^{\max}_{i,t,v}=\max_{x\in B_i(t)}z_v(x),\qquad
\Omega_{i,t,v}=\pi_v(B_i(t)).
\end{equation}
投影足迹向外取整，覆盖全部相交深度像素。单视图要求足迹完全位于图像内、深度全部有效、置信度均不低于 0.8，且不含天空或自车像素，并满足
\begin{equation}
z^{\max}_{i,t,v}<\min_{p\in\Omega_{i,t,v}}
\bigl[D_{t,v}(p)-0.15-0.02D_{t,v}(p)\bigr].
\label{eq:footprint-free}
\end{equation}
满足全部条件的视图投票，其余视图弃权。两票规则、已见表面优先及硬清零机制保持一致。

\begin{table*}[t]
\centering\small
\caption{100 个合成场景的自由空间掩码审计。错误指包含参考表面的单元被标为自由空间，各类别互斥。重建表采用中心射线投票；足迹结果为独立掩码审计。}
\label{tab:free-space-audit}
\begin{tabular}{@{}lrrr@{}}
\toprule
类别 & 含表面单元 & 中心判据误标 & 足迹判据误标 \\
\midrule
薄结构 & 18,437 & 614 & 9 \\
遮挡边界 & 27,619 & 587 & 7 \\
其余区域 & 53,128 & 173 & 3 \\
\midrule
合计 & 99,184 & 1,374 & 19 \\
错误／含表面单元 & & 1.385\% & 0.019\% \\
自由空间标签数量 & & 47,286 & 32,871 \\
含表面污染／自由标签 & & 2.906\% & 0.058\% \\
\bottomrule
\end{tabular}
\end{table*}

足迹判据将占据单元误标从 1,374 个降至 19 个（表~\ref{tab:free-space-audit}）。使用参考深度时，中心判据仍误标 934 个占据单元，将一部分错误定位到单元的横向覆盖。连续参考深度的完整足迹下界在该审计集合中误标为零；离散估计深度下的 19 个剩余错误主要位于亚像素结构和深度边界。
